\bigskip
\centerline{\bf 2. Central Intervals}
\bigskip

\flushpar
Throughout the following sections
we will fix an $S-$unimodal map $f:[-1,1]\to [-1,1]$ with critical exponent 
$\alpha>1$ and without periodic attractors. Furthermore assume that the 
critical orbit is recurrent. 

\flushpar
The set of nice points is
$$
\Cal{N}=\{x\in [-1,1]| \forall i\ge 0 \text{  } f^i(x)\notin (-|x|,|x|)\}
$$ 
This set is closed and not empty. For example the fixed point of $f$ in 
$(0,1)$ is in $\Cal{N}$. 

\flushpar
For $x\in \Cal{N}$ let $D_x\subset U_x=(-|x|,|x|)$ 
be the set of points whose orbit returns to $U_x$. The first return map to 
$U_x$ is denoted by 
$$
R_x:D_x\to U_x.
$$

\flushpar
The next Lemma is a straightforward consequence of
the fact that 
the boundary of each $U_x$ is formed by  nice points.

\proclaim{Lemma 2.1([M])} For every $x\in \Cal{N}$ there exists a collection of
pairwise dispoint intervals $\Cal{U}_x$ with
\parindent=15pt
\item{1)} $I\subset U_x$ for all $I\in \Cal{U}_x$,
\item{2)} $\bigcup_{I\in \Cal{U}_x} I =D_x$,
\item{3)} if $I\in \Cal{U}_x$ and $0\notin I$ then
$R_x:I\to U_x$  
is monotone and onto,
\item{4)} if $I\in \Cal{U}_x$ and $0\in I$ then 
$
R_x:I\to U_x
$
is $2$ to $1$ onto the image. Moreover $R_x(\partial I)= \{x\} 
\text{ or } \{-x\}$.
\endproclaim

\flushpar
Define the function $\psi:\Cal{N}\to\Cal{N}$ by
$$
\psi(x)=\partial V_x\cap(0,1),
$$
where $V_x\in \Cal{U}_x$ is the {\it central} interval: $0\in V_x$. Say
$R_x|_{V_x}=f^{q_x}$ and observe that
$$
\{f(\psi(x), f^2(\psi(x)), \dots, f^{q_x}(\psi(x))=x\}\cap U_x=\emptyset
$$
which follows from the fact that $R_x:D_x\to U_x$ is the first return map.
In particular $\psi(x)\in \Cal{N}$ and we can consider the 
first return map to $V_x$. It will also satisfy Lemma 2.1.

\flushpar
Choose $u_1\in \Cal{N}$
and consider the sequence $u_{n}=\psi(u_{n-1})$ with $n\ge 1$. Use the 
simplified notation $\Cal{U}_n$ for $\Cal{U}_{u_n}$ and denote the first
return maps by
$$
R_n:D_n\to U_n
$$
instead of $R_{u_n}:D_{u_n}\to U_{u_n}$. All these first return maps 
have the  properties stated in Lemma 2.1. Observe that $|U_n|=2u_n$. 

\bigskip

\flushpar
Let $\sigma_n=\frac{u_{n+1}}{u_n}$, $n\ge 1$. We call $\sigma_n$ the 
{\it scaling factor} of level $n$. We will assume that
$$
\sigma_n\to 0.
$$
However, the main Proposition 3.1, can also be proved by using only an 
a priori bound on the scaling factors. The assumption $\sigma_n\to 0$ will
make the exposition less cumbersome.

\proclaim{Lemma 2.2} 
If $I\in \Cal{U}_n$ and $R_n|_I=f^t$ then there exists an interval 
$J\supset f(I)$ such that 
$$
f^{t-1}:J\to U_{n-1}
$$ 
is monotone onto. In particular all the maps $f^{t-1}: f(I) \to U_n$, 
$I\in \Cal{U}_n$ have uniformly bounded distortion. Moreover these maps will
be essentially linear when $n\to \infty$. 
\endproclaim

\demo{Proof} 
Let $I\in \Cal{U}_n$ with $R_n|I=f^t$ and let 
$J\supset f(I)$ be the maximal interval
on which $f^{t-1}$ is monotone. The maximality implies the existence of $i<t-1$
such that $0\in \partial f^i(J)$. Observe that $f^i(f(I))\cap U_n=\emptyset$,
the first return happens after $t-1>i$ steps. 
So $u_n \text{ (or } -u_n\text{)}\in
f^i(J)$. We observed before that the orbit of $f(u_n)$ never enters $U_{n-1}$,
$f^{t-1}(J)\supset U_{n-1}$.
\hfill\hfill\qed $\,\,$ (Lemma 2.2)
\enddemo

\proclaim{Lemma 2.3} For $\epsilon>0$ there exists $n_0\ge 1$ such that 
the hyperbolic length of any $I\in\Cal{U}_n$ is small,
$$
\text{hyp}(I,U_n)\le \epsilon \text{ and also } 
\frac{|f(I)|}{|f(U_n)|}\le \epsilon,
$$
whenever $n\ge n_0$.
\endproclaim

\demo{Proof} Let $I\in \Cal{U}_n$, say $R_n|I=f^t|I$. 
The previous Lemma states the existence of an interval $J\supset f(I)$ 
such that $f^{t-1}:(J,f(I))\to (U_{n-1},U_n)$ is monotone onto. For $n$ large
we see that $U_n$ is a very small middle interval in $U_{n-1}$, it has a very
small hyperbolic length. Because the map $f$ has negative Schwarzian
derivative we get that $f(I)\subset J$ has a very small hyperbolic length.
Observe that $f^{-1}(J)\subset U_n$. Otherwise the orbit of $u_n$ would pass 
through $U_{m-1}$. However this is impossible: we saw before that the orbit 
of $\psi(x)=u_n$ does not cross $U_x=U_{n-1}$.
The Lemma will be proved by pulling back the pair $(J,f(I))$ one step more.
\hfill\hfill\qed $\,\,$ (Lemma 2.3)
\enddemo







